\uuid{g2eO}
\titre{Image réciproque}
\auteur{bodin ; hanani ; mzari}
\organisation{exo7}
\licence{CC BY-NC-SA 4.0}
\datecreate{2020-01-20}
\niveau{L1}
\module{Algèbre}
\chapitre{Logique, ensemble, raisonnement}
\sousChapitre{Ensemble}
\difficulte{1}
\theme{Ensembles, applications, image réciproque}
\exo7qcmid{lille-1.ensembles.009}

\contenu{
\begin{question}
\qclassification{100.02,101.01,102.01,102.02}
\qtags{motcle=image réciproque}
On considère l'application $f:\{1,2,3,4\}\to \{1,2,3,4\}$ définie par
$$f(1)=2,\quad f(2)=3,\quad f(3)=4,\quad f(4)=2.$$
Quelle est la bonne réponse ?
\begin{answers}
    \bad{$f^{-1}(\{2\})=\{1\}$}
    \bad{$f^{-1}(\{2\})=\{3\}$}
    \bad{$f^{-1}(\{2\})=\{4\}$}
    \good{$f^{-1}(\{2\})=\{1,4\}$}
\end{answers}
\begin{explanations}
L'ensemble $f^{-1}(\{2\})$ est formé des éléments qui ont une image égale à $2$.
\end{explanations}
\end{question}
}
